\section{Referee extraction and its clock-grid cost}\label{sec:referee}
Let $G=(X,Y,\mathcal A,\mathcal B,\pi,V)$ be finite. Each party has
a system register, a clock, old question, answer and work registers.
The initialization state of the old registers is
\[
 \ket{\pi,0}=\sum_{x,y}\sqrt{\pi(x,y)}\ket{x,y}\ket{0}_{\rm ans,work}.
\]
The local circuits $U_A,U_B$ read their own question without changing
it; neither uses the other's question. Their coherent outputs are tested
by the diagonal losing projection $\Pi_{\rm lose}$. Put
$Q=(U_A\otimes U_B)^*\Pi_{\rm lose}(U_A\otimes U_B)$.
Equation~\eqref{eq:cell} is an exact assumption on a legal bipartite
tensor product; the two clocks belong to their respective parties.
Extra idle clocks can belong entirely to Alice. Product rewinding is
the controlled local prefix circuit on each party's own clock and data.
No nearest-neighbour realization of this whole initialization projector
is implicit in the assumption.

On fresh external questions drawn with the original joint law $\pi$,
run these circuits on the system with new local pure answer/work
ancillas. All old registers and clocks are spectators. The resulting
loss effect on the old register space is
\begin{equation}\label{eq:loss-effect}
 F=I_{\rm old}\otimes F_{\rm sys},\qquad 0\leq F\leq I,
 \qquad [F,P]=0,\qquad PFP=PQP.
\end{equation}
The last identity follows by expanding the coherent old questions:
preserved labels and the diagonal test kill cross terms, including
when $\pi$ is correlated. In particular
$\bra{Pv}F\ket{Pv}=\norm{QPv}^2$.
Keeping spectators preserves the shared Schmidt data. There is no
postselection at a clock corner.

\begin{lemma}[Elementary two-clock extraction]\label{lem:elementary}
For two clocks, a legal unit vector of energy $\lambda$ yields loss
at most $5(1+\sqrt{2T})^2M\lambda$ under the fresh-question strategy.
If $\eps_{\rm opt}$ is the least initialized loss with these circuits,
$E_0(H)\leq\eps_{\rm opt}/M$.
\end{lemma}
\begin{proof}
Write $W^*\Psi=\sum_v\ket v\psi_v$. Positivity gives
\[
 E_{\rm prop}=\sum_{\{v,w\}}\norm{\psi_v-\psi_w}^2\leq\lambda,
 \quad\norm{(I-P)\psi_0}^2\leq\lambda,
 \quad\norm{Q\psi_t}^2\leq\lambda,
\]
where $t=(T,T)$. A path of length at most $2T$ to either corner gives
$\norm{\psi_v-\psi_c}\leq\sqrt{2T\lambda}$ by Cauchy--Schwarz.
Set $a=1+\sqrt{2T}$. Then
$\norm{(I-P)\psi_v}\leq a\sqrt\lambda$ and
$\norm{QP\psi_v}\leq2a\sqrt\lambda$.
By \eqref{eq:loss-effect}, the cross term between initialized and
uninitialized components vanishes. Summing their losses over the $M$
orthogonal clock components gives $M(4a^2+a^2)\lambda$.
For the variational bound take the uniform clock with an initialized
optimal system input. Propagation and initialization vanish; the final
corner penalty is $\eps_{\rm opt}/M$.
\end{proof}

\begin{lemma}[Hilbert-valued resistance inequality]\label{lem:resistance}
For a connected graph Laplacian $L$, a zero-sum real vector $w$,
and Hilbert-space vectors $\psi_v$,
\begin{equation}\label{eq:resistance}
 \norm{\sum_vw_v\psi_v}^2\leq(w^*L^+w)
          \sum_{\{v,z\}}\norm{\psi_v-\psi_z}^2.
\end{equation}
\end{lemma}
\begin{proof}
Diagonalize $L$ on the orthogonal complement of constants. In that
basis the left vector is the inner product of the scalar coefficients
$L^{+1/2}w$ with the Hilbert-valued coefficients $L^{1/2}\psi$.
Hilbert-space Cauchy--Schwarz gives \eqref{eq:resistance}. The constant
mode drops out because $w$ sums to zero.
\end{proof}

\begin{proof}[Proof of Theorem~\ref{thm:extraction}: energy and price]
Let $m=M^{-1}\sum_v\psi_v$, $d_v=\psi_v-m$, and $G=G_k$.
Using $w=\delta_c-M^{-1}\one$ in Lemma~\ref{lem:resistance} gives
$\norm{d_c}\leq\sqrt{G\lambda}$ at either corner. Poincar\'e's
inequality gives $\sum_v\norm{d_v}^2\leq\lambda/\mu_1$.
The corner initialization and test imply
\begin{align*}
 \norm{(I-P)m}&\leq(1+\sqrt G)\sqrt\lambda,\\
 \norm{QPm}&\leq\sqrt\lambda+\norm{(I-P)m}+2\sqrt{G\lambda}.
\end{align*}
Define $\ell(v)=\bra vF\ket v$. Its square root is a seminorm,
and \eqref{eq:loss-effect} implies
\[
 \sqrt{\ell(m)}\leq\norm{QPm}+\norm{(I-P)m}
                       \leq(3+4\sqrt G)\sqrt\lambda.
\]
Consequently
\begin{equation}\label{eq:extraction-proof}
 \sum_v\ell(\psi_v)\leq2M\ell(m)+2\sum_v\norm{d_v}^2
            \leq C_kM\lambda.
\end{equation}
The shared state is exactly $W^*\Psi$, augmented only by local pure
ancillas. Since $W$ is a product unitary, its Schmidt coefficients,
rank and entropy are unchanged. Compressing both local POVMs to their
Schmidt supports yields a legitimate strategy of local dimension
$\SR(\Psi)$, proving the dimension bound. Pure-state one-way
distillation equals marginal entropy \cite{bbps}; the definition
\eqref{eq:price} proves the entropy bound.

The quadratic-form inequality \eqref{eq:extraction-proof} is an operator
inequality for a fixed fresh-question loss effect. It therefore holds
after taking expectation in any legal density matrix. Local unitary
invariance of $\ED$ and \eqref{eq:price} prove the mixed-state
certificate. No entropy of a mixed marginal is being identified with
its entanglement.
\end{proof}

\subsection{The resistance trichotomy}
The path eigenvalues and corner weights are
\[
 \mu_a=4\sin^2\frac{\pi a}{2n},\quad
 w_0=1/n,\quad w_a=\frac2n\cos^2\frac{\pi a}{2n}\quad(1\leq a<n).
\]
Thus, exactly,
\begin{equation}\label{eq:greensum}
 G_k=\sum_{\substack{0\leq a_1,\ldots,a_k<n\\(a_1,\ldots,a_k)\ne0}}
       \frac{w_{a_1}\cdots w_{a_k}}{\mu_{a_1}+\cdots+\mu_{a_k}}.
\end{equation}
The first path eigenvalue satisfies $\mu_1\geq4/n^2$. For $k\geq2$,
$M\mu_1\geq4n^{k-2}$, proving the variance term in the theorem.

For a path, the centered potential solving
$Lg=\delta_0-n^{-1}\one$ has edge differences
$g_j-g_{j+1}=(n-1-j)/n$. Hence its energy, equal to $G_1$, is
\begin{equation}\label{eq:pathgreen}
 G_1=\frac1{n^2}\sum_{j=1}^{n-1}j^2
       =\frac{T(2T+1)}{6(T+1)}.
\end{equation}
For two paths use $w_a\leq2/n$ and
$\mu_a\geq4a^2/n^2$ to obtain
\[
 G_2\leq\sum_{\substack{0\leq a,b\leq T\\(a,b)\ne(0,0)}}\frac1{a^2+b^2}.
\]
The axes contribute at most $\pi^2/3$ and $(1,1)$ contributes $1/2$.
Each remaining positive-index term is bounded by the integral of
$(x^2+y^2)^{-1}$ over $[a-1,a]\times[b-1,b]$. Those cells lie inside
the quarter annulus $1\leq\sqrt{x^2+y^2}\leq\sqrt2T$. Therefore
\begin{equation}\label{eq:greenbound}
 G_2\leq\frac\pi2\ln T+\frac{\pi^2}3+\frac12+\frac\pi4\ln2
          <\frac\pi2\ln T+4.4\qquad(T\geq2).
\end{equation}
To identify the leading constant, for $a,b\geq1$ compare the exact
summand in \eqref{eq:greensum} with
$4/[\pi^2(a^2+b^2)]$. Taylor bounds for sine and cosine, and
$\sin x\geq2x/\pi$ on $[0,\pi/2]$, bound the difference by
$O(n^{-2})$ uniformly. Summing contributes $O(1)$; the axes are
bounded. The quadrant sum of $(a^2+b^2)^{-1}$ is
$(\pi/2)\ln n+O(1)$ by unit-square integration, proving
$G_2=(2/\pi)\ln T+O(1)$.

For three paths, \eqref{eq:greensum} implies
\[
 G_3\leq\frac2n\sum_{a\ne0}|a|^{-2}
 \leq\frac2n\sum_{m=1}^T\left(3+\frac3m+\frac1{m^2}\right)<14.
\]
The shell $\max_i a_i=m$ has $3m^2+3m+1$ vectors. The last bound
uses $1/m,1/m^2\leq1$ and $T/n<1$. Thus
\[
 C_3\leq2(3+4\sqrt{14})^2+1/2<647.
\]
For any fixed $k>3$ the same shell estimate is $O(T^{k-2})$ before
the prefactor $2^{k-2}n^{2-k}$, so $G_k=O(1)$.
This completes the resistance part of Theorem~\ref{thm:extraction}.
The dimensional phenomenon is classical network theory
\cite{cserti,doyle}; the application is the extraction inequality.

\begin{proposition}[The logarithmic obstruction for all low-energy vectors]\label{prop:lognecessary}
No constant independent of $T$ can replace the two-clock extraction
factor for all legal low-energy vectors in this general cell class.
\end{proposition}
\begin{proof}
Let $c$ be the losing corner, $g=L_2^+\delta_c$, $G=g_c$, and
$f=\one-g/G$. Then $f_c=0$, $M^{-1}\sum_v f_v=1$, and
$f^*L_2f=1/G$. Take $P=I,Q=I$, with the always-losing scalar
system and normalized clock $f/\norm f$. The loss is one and the
energy is $\lambda=1/(G\norm f^2)$, so
\[
 \frac{\text{loss}}{M\lambda}=G\frac{\norm f^2}{M}\geq G.
\]
Since $G\to\infty$, a uniform constant is impossible. A direct-sum
zero-loss system sector can be added. This is an all-low-energy-state
obstruction and makes no assertion about a special game's ground
states alone.
\end{proof}

\subsection{Approximate clock factorization}\label{sec:factorization}
\begin{proposition}[Normalized entropy-support estimate]\label{prop:factorization}
For $k\geq2$ put $u=\sqrt M\,m$ and
$z=\ket{\rm unif}\otimes Pu$. If
$b=(3/2+\sqrt{G_k})\sqrt{M\lambda}<1$, then
\[
 \norm{W^*\Psi-z}\leq b,\qquad
 \norm{W^*\Psi-z/\norm z}\leq2b.
\]
Let $t$ be the actual trace distance between the reduced Alice states
of these two unit vectors. In the small-distance regime,
\begin{equation}\label{eq:entropyupper}
 S_A(\Psi)\leq\log_2d_{{\rm sys},A}+\log_2|X|
                   +t\log_2(D_A-1)+h_2(t),\qquad t\leq2b.
\end{equation}
$D_A$ is an actual local support dimension of the compared states.
\end{proposition}
\begin{proof}
Poincar\'e gives
$\norm{W^*\Psi-\ket{\rm unif}\otimes u}\leq\sqrt{M\lambda}/2$.
Initialization bounds $\norm{(I-P)u}\leq(1+\sqrt G)\sqrt{M\lambda}$.
Their triangle inequality gives the first estimate. Since $\Psi$ is
unit, $|\norm z-1|\leq b$ and $z\ne0$, giving the factor two upon
normalization. The normalized reference consists of product-uniform
local clocks, initialized questions and work, and a system vector.
Its entropy is at most $\log_2|X|+\log_2d_{{\rm sys},A}$.
Local rewinding leaves entropy invariant. The trace distance is no
larger than the vector distance; apply the finite-dimensional
entropy-continuity bound of \cite{audenaert}. One may replace $t$
by $2b$ only where that bound is monotone, for example
$2b\leq1-1/D_A$.
\end{proof}

There is no compulsory additive $\log_2(T+1)$ in this reference
entropy. There is also no exact clock separability at positive energy.
For $P=I,Q=0$, orthogonal path modes $u$ (uniform) and $v$ (first
excited) give the entangled clock state
$(u\otimes u+\delta v\otimes v)/\sqrt{1+\delta^2}$, of energy
$2\mu_1\delta^2/(1+\delta^2)$ and entropy
$h_2(\delta^2/(1+\delta^2))>0$.

For a ground state with $M\lambda\leq\eps_N$, an adequate continuity
condition, if $q$ counts work qubits, is
\[
 \sqrt{\eps_N}(1+\sqrt G)
       [q+\log_2d_{{\rm sys},A}+\log_2(T+1)+O(1)]\longrightarrow0.
\]
It is a condition on the number of qubits, not on its logarithm.

\subsection{Perturbation and dimension certificates}\label{sec:certificates}
\begin{corollary}[Total-norm perturbations]\label{cor:perturbed}
For every legal-space Hermitian $V$ with $\norm V\leq\eta$ and every
pure ground vector $\Psi_\eta$ of $H+V$,
\[
 S_A(\Psi_\eta)\geq\price\bigl(C_k[\eps_{\rm opt}+2M\eta]\bigr),
 \quad \SR(\Psi_\eta)\geq D_G\bigl(C_k[\eps_{\rm opt}+2M\eta]\bigr).
\]
\end{corollary}
\begin{proof}
The variational principle gives
$\bra{\Psi_\eta}H\ket{\Psi_\eta}\leq E_0(H)+2\eta$.
Use Lemma~\ref{lem:elementary} and Theorem~\ref{thm:extraction}.
\end{proof}
Here $\eta$ is the total cell operator norm. Summing local fields may
produce a norm growing with the number of sites.

\begin{proposition}[Register-limited frustration]\label{prop:frustration}
In a fixed classical program block with $n_s$ measured system qubits
per party, let $e$ be the cell ground energy. Then
\[
 D_G(C_kMe)\leq2^{n_s}.
\]
If $\log_2D_G(\eps)=\alpha\log_2(1/\eps)+o(\log(1/\eps))$,
$\alpha>0$, this implies
\begin{equation}\label{eq:frustration}
 e\geq\frac{2^{-n_s/\alpha-o(n_s)}}{C_kM}.
\end{equation}
\end{proposition}
\begin{proof}
After extraction, trace out old registers and clocks locally. The
fresh measurements act only on the measured system, so its mixed
state has the same loss and dimension $2^{n_s}$ on each side. Apply
the definition of $D_G$. For each $\xi>0$ the price is at least
$(\alpha-\xi)\log_2(1/\eps)$ at sufficiently small $\eps$.
Hence $C_kMe\geq2^{-n_s/(\alpha-\xi)}$; outside this small-loss range
the claimed asymptotic lower bound is weaker. Let $\xi$ tend to zero.
The argument is uniform over fixed classical program blocks. A quantum
program used by the measurement must instead be counted in the system
dimension.
\end{proof}

A matching upper frustration law needs a dimension-matched admitted
strategy. Clock propagation does not by itself fix the distinct cell
gap: different almost-winning system sectors can have arbitrarily
close low eigenvalues. Honest uniform histories have energy loss/$M$,
which is a variational scale, not a universal gap formula. A fixed
Hamiltonian minimizes a linear functional of the state; its
measurement encoding must therefore already be specified. Fixed
circuits and classical program blocks suffice, without exhausting
all possible encodings. Finally the elementary extraction factor
shows that $\log T=o(\log(1/\eps))$ is sufficient to preserve a
logarithmic coefficient; it is not a necessary-and-sufficient theorem.

An extension to arbitrary physical states is possible if the legal
projection factorizes across the cut, $W$ has a product-unitary
extension, the illegal space is invariant and has energy at least one,
and the extraction processes its clock labels as spectators. The illegal
weight is then bounded by its energy and its loss by its weight; since
$C_kM\geq1$, the same inequality holds. Block diagonality or a list
of forbidden clock strings alone does not establish these hypotheses.
